\documentclass[12pt]{article} 
\renewcommand{\thempfootnote}{\arabic{mpfootnote}}

\usepackage[hidelinks]{hyperref}

\usepackage[utf8]{inputenc}

\usepackage[T1]{fontenc}
\usepackage[ttdefault=true]{AnonymousPro}

\usepackage{multicol}
\setlength{\columnsep}{1cm}

\usepackage{abc}
\usepackage{graphicx}
\usepackage{musicography}

\usepackage[margin=0.5in]{geometry} 
\geometry{letterpaper}

\usepackage{fancyhdr} 
\pagestyle{fancy} % options: empty , plain , fancy
\renewcommand{\headrulewidth}{0pt} % customise the layout...
\lhead{}\chead{}\rhead{}
\lfoot{}\cfoot{\thepage}\rfoot{}

\begin{document}

\thispagestyle{empty}

\begin{abc}[name=waterfall,options={-O= +c -F single.fmt}]
X:1
T:Over The Waterfall
M:C|
K:D
|: de | "D"f2 a2 "A7"gf e2 | "D"d2 B2 A2 de | "D"f2 a2 "A7"gf e2 | "D" [d4f4] [d2f2] de | 
"D"f2 a2 "A7"gf e2 | "D"d2 B2 A2 AB | "C"=c2 cc B2 A2 | "G"G4 G2 :| 
|: FG | "D"A2A2 "G"B2B2 | "D"ABAG F2 FG | "D"A2 d2 "A7"cd e2 | "D"f4 f4 | 
"D"A2A2 "G"B2B2 | "D"ABAG F2 FG | "D"A2 A2 "A7"GF E2 | "D"D4 D2 :|
\end{abc}
\vskip 0.5in
\noindent\makebox[\textwidth][c]{%
  \begin{minipage}{7.0in}
    \setlength{\parindent}{1.5em}


Over the Waterfall is associated with Virginia fiddler Henry Reed
(1884 - 1968) who believed he heard it from a calliope at a circus or
medicine show when he was 5 or 6 years
old.\footnote{\url{https://tunearch.org/wiki/Annotation:Over\_the\_Waterfall}}
Reed taught it to folklorist and fiddler, Alan Jabbour (1942 – 2017),
who recorded it with the Hollow Rock String Band in
1967.\footnote{\url{https://www.earlyblurs.com/hrsb.htm}} This tune
is believed to have roots in Ireland with its similarities with the
song, Eggs and Marrowbones, and tune, The Dark Girl Dressed in Blue,
but this is difficult to prove.\footnote{tunearch.org, \textit{op. cit.} A similar tune to Over the
  Waterfall can be heard in this recording of Eggs and Marrowbones
  (\url{https://youtu.be/\_veDpbB7lpY}), but I think that The Dark
  Girl Dressed in Blue
  (\url{https://tunearch.org/wiki/Dark\_Girl\_Dressed\_in\_Blue\_(2)\_(The)})
  is much more similar to Over the Waterfall.  These tunes are also
  called Old Woman from Wexford, Old Woman in Dover, and The Fellow
  that Looks Like Me.}

Alan Jabbour is from Jacksonville, FL and was originally trained in
classical violin, even playing with the Jacksonville
Symphony.\footnote{I lived in Jacksonville, FL during high school. It
  was a nice town.} While in graduate school he became interested in
fiddling, collecting Old Time tunes in the Upper South (North
Carolina, Virginia, and West Virginia) with this collection now in the
Archive of Folk Culture at the Library of Congress. His research
eventually led him to head the Archive of Folk Song (now the Archive
of Folk Culture) in 1969, working for the National Endowment for the
Arts in 1974, and perhaps what is he best known for, becoming the
founding director of the American Folklife Center in the Library of
Congress in
1974.\footnote{https://en.wikipedia.org/wiki/Alan\_Jabbour} It is hard
to overstate the influence Alan Jabbour has had on the advocacy
and endorsement of modern American Ethnomusicology.

---Kitty Steetle
  \end{minipage}
}


\end{document}
